\begin{theorem}[“二值跳跃”的统一表述：取向函子同时控制解析跳跃与离散 uplift]\label{thm:xi-binary-jump-orientation-unification}
令 \(\pi:\widetilde Y\to Y\) 为一族局部有限层覆叠，使得对每个 \(y\in Y\) 的纤维 \(S_y:=\pi^{-1}(y)\) 为有限集合；
允许的“重标号”由 \(\Aut(S_y)\cong \mathfrak S_{|S_y|}\) 给出。
对每个纤维定义其取向 \(\ZZ_2\)-torsor
\[
\mathrm{Or}(S_y):=\mathrm{Bij}([|S_y|],S_y)\big/\mathfrak A_{|S_y|}.
\]
则在本文所出现的两类典型二值跳跃机制中，二值对象都可自然表达为 \(\mathrm{Or}\) 的像：
\begin{enumerate}
\normalfont
  \item \textbf{离散 uplift：}对 boundary uplift 位移集合 \(\Delta_m^{\mathrm{bdry}}(w)\)（作为无序有限集），在不引入外加寄存器且要求置换自然性的约束下，唯一的非平凡二值跳跃对象可取为
  \(\mathrm{Or}(\Delta_m^{\mathrm{bdry}}(w))\)（推论 \ref{cor:bdry-orientation-parity-uplift}）。
  \item \textbf{解析跳跃：}在 Joukowsky 二重覆盖的两分支表示中（定理 \ref{thm:xi-joukowsky-cauchy-transform-two-branch}），对椭圆边界点 \(w(\theta)=J_L(e^{\mathrm{i}\theta})\) 存在两侧边界值，
  且 Cauchy 跳跃
  \(
  G_{\sigma_L}^{\mathrm{ext}}(w(\theta))-G_{\sigma_L}^{\mathrm{int}}(w(\theta))
  \)
  在交换两支（等价地在双覆盖纤维 \(\{z_+(w(\theta)),z_-(w(\theta))\}\) 上作用非平凡置换）下变号（定理 \ref{thm:xi-joukowsky-ellipse-plemelj-jump} 的平方根分支互换律）。
  因而该跳跃量可视为取向 torsor \(\mathrm{Or}(\{z_+(w(\theta)),z_-(w(\theta))\})\) 的符号等变截面，从而不依赖对两支的标号选择。
\end{enumerate}
因此，在“允许重标号但禁止外加寄存器”的共同口径下，二值跳跃的可审计结构由取向函子统一刻画；并且在局部层数 \(\ge 3\) 的离散情形中，该二值对象在同构意义下由符号表示唯一确定（定理 \ref{thm:bdry-binary-jump-orientation-functor-uniqueness}）。
\end{theorem}

\begin{proof}
离散部分由推论 \ref{cor:bdry-orientation-parity-uplift} 直接给出。
解析部分仅用对称性：二重覆盖的纤维为二元集合，其非平凡置换生成 \(\mathfrak S_2\)；
而跳跃互换律表明所述差量对该置换取负，故其变换律正是符号表示，从而可在取向 torsor 上无歧义地表达为符号等变截面。
最后一段的“局部层数 \(\ge 3\) 唯一性”由定理 \ref{thm:bdry-binary-jump-orientation-functor-uniqueness} 的分类结论给出。证毕。
\end{proof}

\endinput
